\chapter{Cài đặt hệ thống}

\section{Cài đặt thiết bị nhúng ESP32}

\subsection{Cấu hình I2S}

\begin{table}[htbp]
\centering
\caption{Cấu hình I2S cho microphone INMP441}
\label{tab:i2s-config}
\begin{tabular}{>{\raggedright\arraybackslash}p{0.35\linewidth} >{\raggedright\arraybackslash}p{0.55\linewidth}}
\toprule
\rowcolor{headerblue} \textcolor{white}{\textbf{Tham số}} & \textcolor{white}{\textbf{Giá trị}} \\
\midrule
\rowcolor{rowgray} Chế độ & I2S\_MODE\_MASTER | I2S\_MODE\_RX \\
Tần số lấy mẫu & 16.000 Hz \\
\rowcolor{rowgray} Bit depth đầu vào & 32-bit (I2S native) \\
Bit depth đầu ra & 16-bit PCM (sau chuyển đổi) \\
\rowcolor{rowgray} Kênh & Mono (chỉ kênh trái) \\
Số DMA buffer & 16 (tăng từ 8 để ổn định) \\
\rowcolor{rowgray} Kích thước buffer & 1024 mẫu (64 ms mỗi buffer) \\
Throughput & 32 KB/s ($\approx$ 256 kbps) \\
\bottomrule
\end{tabular}
\end{table}

\subsection{Chuẩn hóa tín hiệu}

Microphone INMP441 có độ phân giải 18-bit, đóng gói trong word 32-bit với dữ liệu
nằm ở các bit cao. Quá trình chuẩn hóa:

\begin{enumerate}
\item Đọc mẫu 32-bit raw từ I2S DMA buffer.
\item Dịch phải 14 bit: \texttt{val = raw\_samples[i] >> 14}.
\item Cắt ngưỡng vào phạm vi 16-bit signed: $[-32768, 32767]$.
\item Lưu dưới dạng \texttt{int16\_t} để truyền qua TCP.
\end{enumerate}

\subsection{Pipeline streaming video}

Camera OV5640 được cấu hình chụp JPEG với double buffering (2 frame buffer).
XCLK mặc định 12 MHz, chất lượng JPEG có thể điều chỉnh (10--63).

Mỗi frame được truyền qua TCP port 3001 với giao thức:
\texttt{[Length (4 bytes LE)][JPEG Data (variable)]}.
Buffer tĩnh 250 KB được cấp phát một lần trên PSRAM (hoặc heap) để tránh
phân mảnh bộ nhớ.

\subsection{Giao thức điều khiển camera}

Trình duyệt gửi lệnh điều khiển qua WebSocket đến server trung gian,
server chuyển tiếp qua TCP đến ESP32 với format:
\texttt{[0xA5][ID (1 byte)][Value (1 byte)]}.

\begin{table}[htbp]
\centering
\caption{Các lệnh điều khiển camera chính}
\label{tab:camera-commands}
\begin{tabular}{>{\centering\arraybackslash}p{0.08\linewidth} >{\raggedright\arraybackslash}p{0.3\linewidth} >{\raggedright\arraybackslash}p{0.45\linewidth}}
\toprule
\rowcolor{headerblue} \textcolor{white}{\textbf{ID}} & \textcolor{white}{\textbf{Tham số}} & \textcolor{white}{\textbf{Phạm vi}} \\
\midrule
\rowcolor{rowgray} 1 & JPEG Quality & 10--63 (thấp = chất lượng cao hơn) \\
2 & Frame Size & QQVGA đến UXGA (14 mức) \\
\rowcolor{rowgray} 3--5 & Brightness, Contrast, Saturation & $-2$ đến $+2$ \\
6--9 & AWB, AEC, AGC, AEC2 & 0/1 (bật/tắt) \\
\rowcolor{rowgray} 16--17 & H-Mirror, V-Flip & 0/1 \\
20 & XCLK Frequency & 6/8/12/16/24 MHz \\
\rowcolor{rowgray} 21 & Frame Buffer Count & 1--4 \\
\bottomrule
\end{tabular}
\end{table}

ESP32 phản hồi cài đặt hiện tại qua packet với magic header \texttt{0xBEEF}
theo format: \texttt{[0xBE][0xEF][ID][Value]...} cho tất cả 21 tham số.

\section{Cài đặt dịch vụ khử nhiễu}

\subsection{Huấn luyện trên VIVOS}

\begin{table}[htbp]
\centering
\caption{Cấu hình huấn luyện DTLN trên VIVOS}
\label{tab:training-config}
\begin{tabular}{>{\raggedright\arraybackslash}p{0.3\linewidth} >{\centering\arraybackslash}p{0.2\linewidth} >{\raggedright\arraybackslash}p{0.35\linewidth}}
\toprule
\rowcolor{headerblue} \textcolor{white}{\textbf{Tham số}} & \textcolor{white}{\textbf{Giá trị}} & \textcolor{white}{\textbf{Giải thích}} \\
\midrule
\rowcolor{rowgray} Dữ liệu sạch & VIVOS & Tiếng Việt, $>$ 1 GB \\
Dữ liệu nhiễu & Microsoft DNS & Kết hợp random SNR \\
\rowcolor{rowgray} Batch size & 32 & Cân bằng GPU/gradient \\
Độ dài đoạn & 3 giây & Phù hợp VIVOS (3--5s) \\
\rowcolor{rowgray} Epoch tối đa & 50 & Kết hợp dừng sớm \\
Learning rate & $10^{-3}$ & ReduceLROnPlateau \\
\rowcolor{rowgray} Early Stopping & 10 epoch & Tránh overfitting \\
Optimizer & Adam, clipnorm=3.0 & Tránh bùng nổ gradient \\
\bottomrule
\end{tabular}
\end{table}

\subsection{Pipeline xử lý khử nhiễu}

Luồng xử lý khi nhận request:
\begin{enumerate}
\item Đọc file WAV qua SoundFile, validate sample rate (16 kHz).
\item Convert stereo→mono (nếu cần), cast sang float32.
\item Thêm batch dimension, chạy \texttt{model.predict()}.
\item Clip output vào $[-1, 1]$, ghi file WAV đã khử nhiễu.
\item Trả file qua HTTP response, xóa file tạm.
\end{enumerate}

\section{Cài đặt dịch vụ nhận dạng giọng nói}

\subsection{Speech-to-Text}

\begin{itemize}
\item \textbf{Mô hình:} OpenAI gpt-4o-transcribe (cloud API).
\item \textbf{Định dạng hỗ trợ:} WAV, MP3, FLAC, OGG, WebM, M4A.
\item \textbf{Ngôn ngữ:} Tiếng Việt (đa vùng miền).
\item \textbf{API:} \texttt{POST /stt} -- upload file audio, trả về text.
\end{itemize}

\subsection{Sinh mô tả sản phẩm}

\begin{itemize}
\item \textbf{Mô hình:} GPT-4o-mini (temperature 0.75, max 600 tokens).
\item \textbf{API:} \texttt{POST /gen} -- nhận \texttt{product\_description} (text từ STT)
  và \texttt{style} (phong cách viết).
\item \textbf{Prompt engineering:} System prompt hướng dẫn AI đóng vai copywriter
  thương mại điện tử, trích xuất thông tin sản phẩm từ text lộn xộn của STT,
  bỏ qua từ lấp, lặp lại, lỗi phát âm; viết 3--5 câu ngắn gọn theo phong cách yêu cầu.
\end{itemize}

\subsection{Lưu trữ phong cách viết}

Các phong cách viết được lưu trong SQLite (\texttt{styles.db}) với schema đơn giản
(id, name, description). Người dùng có thể thêm phong cách mới qua
\texttt{POST /add-style}.

\section{Cài đặt dịch vụ phân loại sản phẩm}

\subsection{Pipeline suy luận}

\begin{enumerate}
\item Client gửi text mô tả sản phẩm qua POST request.
\item Validate input bằng Pydantic schema.
\item Tokenize input (max 512 tokens) qua AutoTokenizer.
\item Forward pass sinh logits cho 32 danh mục.
\item Softmax → top-K predictions với category ID, name, display name, score.
\item Trả response kèm inference time.
\end{enumerate}

Model được load khi khởi động qua FastAPI Lifespan Context Manager,
tự động detect GPU (fallback CPU), set evaluation mode.

\subsection{Hiệu năng}

\begin{table}[htbp]
\centering
\caption{Benchmark hiệu năng phân loại sản phẩm}
\label{tab:classify-benchmark}
\begin{tabular}{>{\raggedright\arraybackslash}p{0.35\linewidth} >{\centering\arraybackslash}p{0.2\linewidth} >{\centering\arraybackslash}p{0.2\linewidth}}
\toprule
\rowcolor{headerblue} \textcolor{white}{\textbf{Scenario}} & \textcolor{white}{\textbf{Latency}} & \textcolor{white}{\textbf{Throughput}} \\
\midrule
\rowcolor{rowgray} Single inference (CPU) & $\approx$ 45 ms & 22 req/s \\
Batch 10 items (CPU) & $\approx$ 180 ms & 55 items/s \\
\rowcolor{rowgray} Single inference (GPU) & $\approx$ 15 ms & 66 req/s \\
Batch 100 items (GPU) & $\approx$ 800 ms & 125 items/s \\
\bottomrule
\end{tabular}
\end{table}
